\documentclass[12pt]{article}
\usepackage[utf8]{inputenc}
\usepackage[T1]{fontenc}
\usepackage{hyperref}
\title{Humla Language Resources for K-12 Curriculum Development}
\author{Compiled by Claude Sonnet 5.5 (claude-sonnet-5-5) for Andr\'as Kornai}
\date{}
\begin{document}
\maketitle

\section{Introduction}
Humla (ISO 639-3 hut), also called Humla Tibetan, is spoken in Humla district, Nepal, by about 5,000 people, with Upper, Lower, Eastern and Limi dialects. Searches found no grammar, dictionary or school materials.

\section{Primary Educational Resources}
None found in search results.

\section{Linguistic and Reference Materials}
\begin{itemize}
\item Humla Tibetan language (Wikipedia). \url{https://en.wikipedia.org/wiki/Humla_Tibetan_language}
\item Klaas de Vries (2014), \emph{A Sociolinguistic Survey of Humla Tibetan, Northwest Nepal} (Glottolog reference record). \url{https://glottolog.org/resource/reference/id/552808}
\item Christopher Wilde (2001), MA thesis on Limi dialect phonology, University of Helsinki (ASJP source record). \url{https://asjp.clld.org/sources/3045}
\end{itemize}

\section{Religious and Community Resources}
Joshua Project's page for Humla lists a Bible portion but no complete Bible. \url{https://www.joshuaproject.net/index.php/languages/hut}

\section{Digital and Audio Resources}
None found in search results.

\section{Curriculum Development Context}
Very little exists; a curriculum would need new documentation work with the community, building on the phonology and survey sources above.

\end{document}